\subsection{ACT-ViT}
\label{sec:act_vit_subsec}

\paragraph{Overview.} ACT-ViT consists of the following components (see Figure~\ref{fig:overview}): (1) a Pooling (Pool) layer
that reduces the ``spatial'' dimensions of the activation tensor (layers and tokens) to a fixed, predefined
size $(L_p, N_p)$; (2) a per-LLM Linear Adapter (LA) layer that compresses and ``aligns'' their feature
dimensions; (3) a ViT-Based Backbone (ViT-B), which processes the resulting tensor using a Vision
Transformer architecture, enhanced with shared positional encodings within each patch. In what
follows, we expand on each of these components.

\paragraph{Pool.} Similarly to image resizing, we propose to reduce the size
of activation tensors before processing them with our architecture: as a preprocessing step, we apply a max-pooling operation
across the spatial activation dimensions (see inset illustration).
Specifically, we propose two possible pooling approaches. The
first maintains full input resolution at the level of tokens and
applies 1d pooling on the layer dimension only. This approach
is preferred when expecting relatively concise LLM generations, or when their maximum generation
length remains computationally manageable. Alternatively, a second approach adopts a more general
2d pooling on both the two spatial activation dimensions, allowing to more efficiently deal with
potentially longer LLM responses. As we will detail next, we will mainly adopt the first approach,
and additionally experiment with the second in Section~\ref{sec:experiments}.

Note that, in our applications, not only the number of generated tokens $N$ naturally varies depending on the input prompt, also, the number of LLM layers $L$ may vary across models in $\mathcal{M}$.
With guiding principle (i) in mind, it becomes important to ensure our architecture can robustly
deal with both these varying parameters. Accordingly, we design the two above pooling operations in a way that, regardless of the input size, their outputs, $A_p \in \mathbb{R}^{L_p \times N_p \times D_M}$, always match
a predefined shape of $(L_p, N_p)$ in their spatial dimensions (see Algorithm~\ref{alg:pool}, in Appendix~\ref{app:pool}). Importantly, this property ensures the downstream ViT backbone is always fed with tensors of a consistent spatial shape.

\paragraph{LA.} Following pooling, we map the feature dimensions to
a shared hidden size via a linear transformation. In particular, we propose to instantiate and employ a dedicated linear
adapter (LA) module for each LLM $M \in \mathcal{M}$; see inset illustration. All adapters project to a shared output dimensionality
$D'$. Formally:
\begin{equation}
\mathcal{L} = \left\{ LA_M : A_p \mapsto A_p W_M \mid W_M \in \mathbb{R}^{D_M \times D'},\ M \in \mathcal{M} \right\},
\end{equation}
where $D_M$ is the hidden dimension of model $M$, and $A_p$ denotes the pooled activation tensor. The
result is a tensor of shape $L_p \times N_p \times D'$. This design---applying a single linear transformation, $LA_M$,
jointly across all layers and tokens of the activation tensor for a given LLM $M$, rather than using
separate projections per layer or token---is motivated by the inherent alignment of feature spaces
within the architecture of a single transformer model. Specifically, residual connections promote
consistency across layers, while the weight sharing of the transformer across tokens ensures a uniform
representation structure along the token dimension.

Resorting to a set of LLM-dependent LAs is particularly relevant in relation to principle (i). Indeed,
this expedient ensures that the downstream architectural components process objects which are
of consistent shape and are meaningfully comparable across the internal representation spaces
of the distinct LLMs in $\mathcal{M}$. Indeed, motivated by recent literature on representation alignment
\citep{huh2024,morcos2018,moschella2022,kornblith2019}, we hypothesize that, despite differences in architectural design, training corpora,
weight initializations and learning schemes, the hidden representations of LLMs can, in fact, be
put in an approximate correspondence by some linear transformations, which we aim to learn via
the adapters in $\mathcal{L}$. This contributes to rendering our multi-LLM setup not only computationally
possible, but also meaningful from a learning perspective and, as we show in the next sections, also
beneficial in terms of generalization performance. For an illustrative motivating example of how a
linear mapping can tackle misalignment in the representation space, we refer the interested reader
to Appendix~\ref{app:permutation}, where we specifically consider those induced by neuron symmetries, as recently
studied in \citep{navon2023,samuelkainsworth2022}.

\paragraph{ViT-B.} In view of the proposed analogy between ATs and images, the output of any LA module is
then processed \`a la ViT \citep{dosovitskiy2020}. First, it is rearranged into a sequence of non-overlapping activation
patches of size $(p_H, p_W)$. If $L_p$ and $N_p$ are divisible by, resp., $p_H$ and $p_W$, this gives $H = L_p/p_H$ and
$W = N_p/p_W$ patches. The resulting reshaped tensor is denoted: $A_{\text{patches}} \in \mathbb{R}^{(H \cdot W) \times p_H \times p_W \times D'}$.
We augment each patch with a per-patch learnable Positional Encoding (PE), providing intra-patch
ordering over (pooled) activation pixels. The Transformer's native PEs provide a global ordering
over activation patches. Combined with our per-patch PEs, this allows reconstructing each pooled
pixel's original position. Similarly to ViT, the enriched activation patches are flattened and passed
through a shared linear layer. The resulting sequence is then processed by a stack of Transformer
blocks. Overall, using a vision-like backbone aligns with the aforementioned analogy between ATs
and images and, as we will experimentally show next, it reveals to be an effective inductive bias.
